\chapter{Open Doors in Algebra}\label{conjectures:ch:algebra}

\begin{center}\itshape
In algebra the hardest questions are often the ones a first-year student can state.\\
\upshape(our summary of this chapter)
\end{center}

Galois taught us to attach a finite group to every polynomial. Tome~2 computed these groups for small examples, Tome~12 read them off the
factorisations of the polynomial modulo primes, and in every case the question was: \emph{here is a polynomial, what is its group?} In 1892 David Hilbert
turned the question around. Here is a finite group $G$: is there a polynomial with rational coefficients whose Galois group is $G$? He answered it for the
symmetric and alternating groups. More than a century later, nobody knows the answer in general.

This last chapter visits four questions of algebra that are easy to state and still open, or were open until recently. The \emph{inverse Galois problem}
asks which groups occur over $\Q$. \emph{Lehmer's problem} asks how close to $1$ the product of the large roots of an integer polynomial can be without being
$1$. The \emph{Jacobian conjecture} asks whether a polynomial map with constant non-zero Jacobian determinant always has a polynomial inverse. And a
theorem of Ax and Grothendieck, which \emph{is} proved, says that an injective polynomial map of $\C^n$ to itself is surjective; its proof over finite
fields takes three lines, and it is the right tool for the quadratic case of the Jacobian conjecture.

We prove what is accessible: every finite abelian group and every symmetric group occur over $\Q$, with explicit polynomials; a non-reciprocal polynomial
has Mahler measure at least $1.28$, so that Lehmer's problem is a problem about reciprocal polynomials; a quadratic map with non-vanishing Jacobian is
injective; and injective polynomial maps over algebraic extensions of $\mathbb F_p$ are bijective.

\begin{think}
The polynomial $X^2-2$ has Galois group $\Z/2\Z$, and $X^3-2$ has Galois group $S_3$. Can you find a polynomial with Galois group $\Z/3\Z$? With group
$\Z/2\Z\times\Z/2\Z$? Which roots of unity would you try?
\end{think}

\section{The inverse Galois problem}

Hilbert's tool was his \emph{irreducibility theorem} (1892): if a polynomial $f(t,X)$ with coefficients in $\Q(t)$ has Galois group $G$ over $\Q(t)$, then
for infinitely many rational values $t_0$ of $t$ the specialised polynomial $f(t_0,X)$ still has Galois group $G$ over $\Q$. The general polynomial
$X^n+t_1X^{n-1}+\dots+t_n$ has group $S_n$ over $\Q(t_1,\dots,t_n)$, and specialising one variable at a time gives polynomials over $\Q$ with group $S_n$.
In 1918 Emmy Noether proposed to treat every group the same way: let $G\subseteq S_n$ act on $\Q(t_1,\dots,t_n)$ by permuting the variables, and hope that
the fixed field is again a field of rational functions. When it is, Hilbert's theorem produces $G$ over $\Q$. In 1969 Richard Swan showed that the hope
fails for the cyclic group of order $47$.

The problem then became a hunt, group by group. Igor Shafarevich proved in 1954 that every finite solvable group is a Galois group over $\Q$ (a gap at the
prime $2$ was found and repaired much later). In the 1980s, the \emph{rigidity method} of Belyi, Fried, Matzat and Thompson realised most of the sporadic
simple groups, including John Thompson's realisation of the Monster, of order about $8\cdot10^{53}$, in 1984. By the end of the decade $25$ of the
$26$ sporadic groups were done. The last one, the Mathieu group $M_{23}$, resisted for almost forty years, until a 2026 preprint of Huang, Jackson, Lee,
Poonen, Pries and Zhang gave an explicit polynomial of degree $23$ with group $M_{23}$. The general problem is open: many families of simple groups of Lie
type are still missing.

\subsection*{Abelian groups}

The first family that Galois theory hands us for free is abelian.

\begin{theorem}\label{conjectures:thm:abelian}
Every finite abelian group is the Galois group of a Galois extension of $\Q$ contained in a cyclotomic field $\Q(\zeta_N)$.
\end{theorem}
\begin{proof}
Write $A\cong\Z/n_1\Z\times\dots\times\Z/n_k\Z$ (Tome~2, Corollary~7.7). There are infinitely many primes $\equiv1\pmod n$ for every $n$ (Tome~1,
Exercise~5.7; or Dirichlet's theorem, Tome~5, Theorem~4.1), so we may choose distinct primes $p_1,\dots,p_k$ with $p_i\equiv1\pmod{n_i}$. Let
$N=p_1\cdots p_k$. By the Chinese remainder theorem and the cyclicity of $(\Z/p\Z)^*$ (Tome~1, Theorem~5.3),
\[
(\Z/N\Z)^*\cong\prod_{i=1}^k(\Z/p_i\Z)^*\cong\prod_{i=1}^k\Z/(p_i-1)\Z.
\]
Since $n_i\mid p_i-1$, reduction modulo $n_i$ is a surjective homomorphism $\Z/(p_i-1)\Z\to\Z/n_i\Z$, and the product of these maps is a surjection
$\pi:(\Z/N\Z)^*\to A$. Let $H=\ker\pi$. Now $\operatorname{Gal}(\Q(\zeta_N)/\Q)\cong(\Z/N\Z)^*$ (Tome~2, Exercise~10.8), and every subgroup of an abelian
group is normal. By the fundamental theorem of Galois theory (Tome~2, Theorem~10.6), the fixed field $K$ of $H$ is a Galois extension of $\Q$ with group
$(\Z/N\Z)^*/H\cong A$.
\end{proof}

Exercise~\ref{conjectures:ex:period} makes this explicit for $\Z/5\Z$. The Kronecker--Weber theorem (Tome~12, Theorem~10.1) says that the construction is the only one:
every abelian extension of $\Q$ lies in some $\Q(\zeta_N)$.

\subsection*{Symmetric groups, by the Chinese remainder theorem}

For $S_n$ we do not need Hilbert's irreducibility theorem: Dedekind's theorem (Tome~12, Theorem~2.20) lets us read cycle types off factorisations modulo
primes, and the Chinese remainder theorem lets us choose the factorisations.

\begin{lemma}\label{conjectures:lem:sngen}
Let $n\ge2$, and let $G\subseteq S_n$ be a transitive subgroup that contains an $(n-1)$-cycle and a transposition. Then $G=S_n$.
\end{lemma}
\begin{proof}
Number the points so that the $(n-1)$-cycle $c$ fixes $n$. The powers of $c$ already move $1$ to any of $1,\dots,n-1$; together with transitivity, this
makes $G$ \emph{doubly transitive}: for any two ordered pairs $(a,b)$ and $(i,j)$ of distinct points there is $g\in G$ with $g(a)=i$ and $g(b)=j$. Indeed,
choose $u\in G$ with $u(b)=n$ (transitivity), then a power $c^r$ with $c^r(u(a))=v(i)$ where $v\in G$ is chosen with $v(j)=n$; since $u(a)\ne n$ and
$v(i)\ne n$, such an $r$ exists. Then $g=v^{-1}c^ru$ sends $a$ to $i$ and $b$ to $v^{-1}(c^r(n))=v^{-1}(n)=j$. Now let $\tau=(a\ b)\in G$ be the
transposition. For each pair $i\ne j$ take $g$ as above: $g\tau g^{-1}=(g(a)\ g(b))=(i\ j)\in G$. So $G$ contains every transposition, and the
transpositions generate $S_n$ (Tome~2, Proposition~4.5).
\end{proof}

\begin{theorem}\label{conjectures:thm:sn}
For every $n\ge2$ there is a monic polynomial $f\in\Z[X]$ of degree $n$ whose Galois group over $\Q$ is $S_n$.
\end{theorem}
\begin{proof}
Over every finite field $\mathbb F_p$ there are monic irreducible polynomials of every degree $d\ge1$: the minimal polynomial over $\mathbb F_p$ of a generator
of the cyclic group $\mathbb F_{p^d}^*$ (Tome~2, Theorems~9.4 and~9.6) has degree $d$. Choose monic polynomials of degree $n$:
\begin{itemize}
\item $f_2\in\mathbb F_2[X]$ irreducible;
\item $f_3=X\,g_3\in\mathbb F_3[X]$ with $g_3$ irreducible of degree $n-1$ (for $n=2$, take $f_3=X(X+1)$);
\item $f_5=q\,r\in\mathbb F_5[X]$ with $q$ irreducible of degree $2$ and $r$ a product of distinct monic irreducible polynomials of \emph{odd} degrees, of
total degree $n-2$, none equal to $q$: $r=1$ if $n=2$; $r$ irreducible if $n-2$ is odd; $r=X(X+1)$ if $n=4$; $r=X\,s$ with $s$ irreducible of degree
$n-3\ge3$ if $n-2\ge4$ is even.
\end{itemize}
Each $f_p$ is squarefree, because its irreducible factors are distinct. Lift $f_2,f_3,f_5$ to monic polynomials in $\Z[X]$ with the same names and put
\[
f=-15f_2+10f_3+6f_5.
\]
Since $-15\equiv1\pmod2$, $10\equiv1\pmod3$, $6\equiv1\pmod5$, and the other two coefficients vanish modulo each of these primes, $f\equiv f_p\pmod p$ for
$p=2,3,5$; and $f$ is monic of degree $n$ because $-15+10+6=1$.

\emph{$f$ is irreducible over $\Q$.} By Gauss's lemma a factorisation over $\Q$ gives one into monic factors in $\Z[X]$ of smaller degrees, and reducing it
modulo $2$ would factor the irreducible $f_2$.

\emph{The primes $2,3,5$ do not divide $\operatorname{disc}(f)$.} The discriminant is a polynomial in the coefficients with integer coefficients, so
$\operatorname{disc}(f)\bmod p=\operatorname{disc}(f_p)$, which is non-zero because $f_p$ is squarefree.

\emph{The cycle types.} Let $G\subseteq S_n$ be the Galois group, acting on the $n$ roots; $G$ is transitive because $f$ is irreducible. By Dedekind's
theorem, $G$ contains a permutation of cycle type $(1,n-1)$, read off modulo $3$, which is an $(n-1)$-cycle, and a permutation $\sigma$ of cycle type
$(2,d_1,\dots,d_s)$ with all $d_j$ odd, read off modulo $5$. Let $m=d_1\cdots d_s$, which is odd. Then $\sigma^m$ is the $m$-th power of a
transposition, that is the transposition itself, times the $m$-th powers of cycles whose lengths divide $m$, which are trivial. So $G$ contains a
transposition, and $G=S_n$ by Lemma~\ref{conjectures:lem:sngen}.
\end{proof}

Only the residues of the coefficients modulo $30$ matter, so we may take them between $-15$ and $15$.

\begin{example}\label{conjectures:exa:sextic}
For $n=6$, take $f_2=X^6+X+1$, $f_3=X(X^5-X+1)$ and $f_5=(X^2+2)\,X\,(X^3+X+1)$. Reducing the coefficients of $-15f_2+10f_3+6f_5$ modulo $30$ gives
\[
f=X^6-12X^4+6X^3+2X^2+7X-15,\qquad \operatorname{disc}(f)=951\,221\cdot6\,107\,273,
\]
with Galois group $S_6$. A numerical calculation builds such polynomials for $n\le10$ and checks every factorisation.
\end{example}

\begin{remark}[most polynomials have group $S_n$]\label{conjectures:rem:vdw}
In 1936 van der Waerden showed that, among the monic polynomials of degree $n$ with integer coefficients bounded by $H$, the proportion whose Galois group is
not $S_n$ tends to $0$ as $H\to\infty$, and he conjectured that the exceptions number $O(H^{n-1})$, the count of the reducible ones. Bhargava proved this
conjecture (Annals of Mathematics, 2025). Explicit families are also known: Osada (1987) showed that $X^n-X-1$ has group $S_n$ for every $n$.
\end{remark}

\begin{remark}[what is known]\label{conjectures:rem:igpknown}
Every finite solvable group (Shafarevich), every symmetric and alternating group (Hilbert), every sporadic simple group (with $M_{23}$ in 2026), and many
groups of Lie type are Galois groups over $\Q$; Malle and Matzat's \emph{Inverse Galois Theory} is the reference. Most of these realisations are
\emph{regular}: they come from extensions of $\Q(t)$ and Hilbert's theorem, which gives infinitely many extensions of $\Q$ at once. Whether every finite
group occurs, over $\Q$ or regularly over $\Q(t)$, is open.
\end{remark}

\section{Lehmer's problem}

The \emph{Mahler measure} of a polynomial $P(X)=a\prod_{i=1}^d(X-\alpha_i)\in\C[X]$, $a\ne0$, is
\[
M(P)=|a|\prod_{i=1}^d\max\bigl(1,|\alpha_i|\bigr).
\]
It is multiplicative, $M(PQ)=M(P)M(Q)$, and Jensen's formula (P1, Problem~2.18) writes $\log M(P)$ as the average of $\log|P|$ on the unit circle. For
$P\in\Z[X]$, $M(P)\ge1$, and Kronecker's theorem (Tome~12, Exercise~1.6) says that a monic irreducible $P\in\Z[X]$ with $M(P)=1$ is $X$ or a
cyclotomic polynomial.

In 1933 Derrick Henry Lehmer was looking for large primes among the numbers $\prod_i(\alpha_i^k-1)$, which generalise the Mersenne numbers $2^k-1$. These
numbers grow like $M(P)^k$, so the polynomials of small measure give the slowest growth and the best chance of primes. He asked whether $M(P)$ can be
made arbitrarily close to $1$ without being $1$, and reported that the smallest value he could find was that of
\[
L(X)=X^{10}+X^9-X^7-X^6-X^5-X^4-X^3+X+1,\qquad M(L)=1.176280818\ldots
\]
Ninety years and many computer searches later (among the largest, those of Mossinghoff, Rhin and Wu in 2008), nobody has found a smaller value.

\begin{center}
\emph{Lehmer's problem.} Is there $c>1$ such that every $P\in\Z[X]$ with $M(P)>1$ has $M(P)\ge c$? Is $c=M(L)$?
\end{center}

\begin{figure}[ht]
\centering
\fig{conjectures-ch12-lehmer}
\caption{The roots of Lehmer's polynomial (blue): eight on the unit circle, one real root $1.17628$ outside and its inverse inside. The roots of $x^3-x-1$
(red): the plastic number $1.32472$ and two complex roots inside the circle, not symmetric about it. Computed for this book.}
\label{conjectures:fig:lehmer}
\end{figure}

Figure~\ref{conjectures:fig:lehmer} shows the shape of Lehmer's polynomial: it is \emph{reciprocal}, its coefficients read the same backwards, so its roots are
symmetric under $z\mapsto1/z$. For a polynomial $P$ of degree $d$ write $P^*(X)=X^dP(1/X)$, the polynomial with the coefficients reversed; $P$ is reciprocal
if $P^*=\pm P$. The next theorem shows that the non-reciprocal polynomials stay far from $1$, so that Lehmer's problem is about reciprocal ones.

\begin{lemma}\label{conjectures:lem:coef}
Let $\varphi(z)=\sum_{i\ge0}\gamma_iz^i$ be holomorphic on an open set containing the closed unit disc, with $|\varphi|\le1$ on the unit circle. Then
$|\gamma_k|\le1-|\gamma_0|^2$ for every $k\ge1$.
\end{lemma}
\begin{proof}
By the maximum principle (Tome~4, Theorem~9.19), $|\varphi|\le1$ on the closed disc. If $|\gamma_0|=1$: the series converges uniformly on the circle, so
$\sum_i|\gamma_i|^2=\frac1{2\pi}\int_0^{2\pi}|\varphi(e^{i\theta})|^2\dd\theta\le1$, and $\gamma_k=0$ for $k\ge1$. Let $a=\gamma_0$ with $|a|<1$.

\emph{The case $k=1$.} The denominator of $\psi=\frac{\varphi-a}{1-\bar a\varphi}$ does not vanish on the closed disc, hence on a neighbourhood of it, so
$\psi$ is holomorphic there, with $\psi(0)=0$. For $|w|\le1$, $|w-a|^2-|1-\bar aw|^2=-(1-|w|^2)(1-|a|^2)\le0$, so $|\psi|\le1$ on the circle. Then
$\psi(z)/z$ is holomorphic (Tome~7, Chapter~9) and bounded by $1$ on the circle, hence at $0$ by the maximum principle: $|\psi'(0)|\le1$. Since
$\psi'=\frac{\varphi'(1-|a|^2)}{(1-\bar a\varphi)^2}$, $\psi'(0)=\frac{\gamma_1}{1-|a|^2}$, and $|\gamma_1|\le1-|a|^2$.

\emph{General $k$.} Let $\omega=e^{2i\pi/k}$. Since $\sum_{j=0}^{k-1}\omega^{ji}$ is $k$ if $k\mid i$ and $0$ otherwise,
\[
\frac1k\sum_{j=0}^{k-1}\varphi(\omega^jz)=\sum_{m\ge0}\gamma_{mk}z^{mk}=\Phi(z^k),\qquad\Phi(w)=\sum_{m\ge0}\gamma_{mk}w^m.
\]
The series $\Phi$ converges on a disc of radius $>1$, and on $|w|=1$, writing $w=z^k$ with $|z|=1$, $|\Phi(w)|\le\frac1k\sum_j|\varphi(\omega^jz)|\le1$.
The case $k=1$ applied to $\Phi$ gives $|\gamma_k|\le1-|\gamma_0|^2$.
\end{proof}

\begin{theorem}[after Breusch and Smyth]\label{conjectures:thm:smyth}
Let $P\in\Z[X]$ be monic, irreducible and not reciprocal. Then
\[
M(P)\ge\frac{1+\sqrt{17}}4=1.2807764\ldots
\]
\end{theorem}
\begin{proof}
Let $d=\deg P$ and $\alpha_1,\dots,\alpha_d$ the roots. As $P$ is irreducible and not reciprocal, $P\ne X$, so $P(0)\ne0$, and $|P(0)|=\prod|\alpha_i|\le
M(P)$. If $|P(0)|\ge2$ we are done; so let $P(0)=\pm1$.

\emph{No root lies on the unit circle.} If $|\alpha|=1$, then $1/\alpha=\bar\alpha$ is also a root ($P$ has real coefficients), so $P$ and $P^*$, whose roots are
the $1/\alpha_i$, share a root. Both are irreducible of degree $d$ ($P^*$ is irreducible since a factorisation of $P^*$ reverses to one of $P$), so
$P^*=cP$ with $c$ the leading coefficient $P(0)=\pm1$, and $P$ would be reciprocal.

\emph{A series with integer coefficients.} $P^*(X)=\prod_i(1-\alpha_iX)$ has constant term $1$, so $1/P^*$ is a power series with integer coefficients
(solve $P^*\cdot\sum e_iX^i=1$ term by term, dividing only by $1$), and
\[
G(z)=\frac{P(z)}{P^*(z)}=\sum_{i\ge0}q_iz^i,\qquad q_i\in\Z,\quad q_0=P(0)=\pm1.
\]
$G$ is not constant, otherwise $P=q_0P^*$ would be reciprocal. Let $k\ge1$ be the least index with $q_k\ne0$; then $|q_k|\ge1$.

\emph{Splitting $G$.} Put
\[
g(z)=\prod_{|\alpha_i|<1}\frac{z-\alpha_i}{1-\alpha_iz},\qquad h(z)=\prod_{|\alpha_i|>1}\frac{1-\alpha_iz}{z-\alpha_i},
\]
so that $g=hG$. The poles of $g$ are the $1/\alpha_i$ with $|\alpha_i|<1$ and those of $h$ are the $\alpha_i$ with $|\alpha_i|>1$, all outside the closed
disc. The roots inside the circle are stable under complex conjugation, and so are those outside; so on $|z|=1$,
$\prod_{|\alpha_i|<1}|1-\alpha_iz|=\prod_{|\alpha_i|<1}|1-\bar\alpha_iz|=\prod_{|\alpha_i|<1}|z-\alpha_i|$, using $|1-\bar\alpha z|=|\bar z-\bar\alpha|=
|z-\alpha|$; hence $|g|=1$, and likewise $|h|=1$ on the circle.

\emph{Comparing coefficients.} Write $g=\sum b_iz^i$ and $h=\sum c_iz^i$. Then $|c_0|=|h(0)|=\prod_{|\alpha_i|>1}|\alpha_i|^{-1}=M(P)^{-1}$, and from
$g=hG$ with $q_1=\dots=q_{k-1}=0$:
\[
b_0=c_0q_0,\qquad b_k=c_0q_k+c_kq_0.
\]
So $|b_0|=|c_0|$, and by Lemma~\ref{conjectures:lem:coef} applied to $g$ and to $h$,
\[
\frac1{M(P)}\le|c_0||q_k|=|b_k-c_kq_0|\le|b_k|+|c_k|\le2\Bigl(1-\frac1{M(P)^2}\Bigr).
\]
With $M=M(P)$ this reads $2M^2-M-2\ge0$, that is $M\ge\frac{1+\sqrt{17}}4$.
\end{proof}

\begin{corollary}\label{conjectures:cor:nonrec}
Let $P\in\Z[X]$ with $P(0)\ne0$ be non-reciprocal. Then $M(P)\ge\frac{1+\sqrt{17}}4$.
\end{corollary}
\begin{proof}
Factor $P=\pm\prod P_j$ into irreducible polynomials of $\Z[X]$. If every $P_j$ satisfied $P_j^*=\pm P_j$, then $P^*=\prod(\pm P_j^*)=\pm P$, since
$(QR)^*=Q^*R^*$. So some $P_j$ is not reciprocal; $P_j(0)\ne0$. If its leading coefficient is $\pm1$, Theorem~\ref{conjectures:thm:smyth} applies to $\pm P_j$; otherwise
$M(P_j)\ge2$. In both cases $M(P)\ge M(P_j)\ge\frac{1+\sqrt{17}}4$, since the other factors have measure at least $1$.
\end{proof}

\begin{figure}[ht]
\centering
\fig{conjectures-ch12-mahler}
\caption{For each degree $d\le14$, the smallest Mahler measure $>1$ among the $2\cdot3^{d-1}$ polynomials of degree $d$ with coefficients in
$\{-1,0,1\}$, leading coefficient $1$ and constant term $\pm1$: reciprocal (blue) and non-reciprocal (red). The red points never go below Smyth's bound,
attained by $x^3-x-1$ in degree $3$ and, in higher degrees, by its relatives such as $x^5+x+1=(x^2+x+1)(x^3-x^2+1)$. The blue points reach Lehmer's polynomial at $d=10$; at $d=8$ the reciprocal
$x^8-x^5-x^4-x^3+1$, of measure $1.28064$, sits just below $\frac{1+\sqrt{17}}4$. Exhaustive search for this book.}
\label{conjectures:fig:mahler}
\end{figure}

\begin{remark}[Smyth's sharp bound]\label{conjectures:rem:smythsharp}
In 1951 Breusch proved a bound of this kind, with the constant $1.1796$. In 1971 Chris Smyth found the best constant: a non-reciprocal $P\in\Z[X]$ with $P(0)\ne0$
has $M(P)\ge\theta_0=1.3247179\ldots$, the real root of $x^3-x-1$ (the plastic number), and the extremal polynomials are known. His proof starts as ours,
with $g=hG$, but draws more from Schur's description of the coefficients of bounded holomorphic functions than the single inequality of
Lemma~\ref{conjectures:lem:coef}. Figure~\ref{conjectures:fig:mahler} shows how sharp it is.
\end{remark}

\begin{remark}[the best general bound]\label{conjectures:rem:dobrowolski}
For reciprocal polynomials the best unconditional result is Dobrowolski's (1979): if $\alpha$ is an algebraic integer of degree $d$ which is not a root of
unity, then $M(\alpha)\ge1+c\bigl(\frac{\log\log d}{\log d}\bigr)^3$ for an absolute constant $c>0$. It tends to $1$ as $d$ grows, and nothing better than
this shape is known. The idea is Fermat's little theorem for polynomials, $P(X^p)\equiv P(X)^p\pmod p$: it makes the integer resultant of $P(X)$ and $P(X^p)$,
a product of differences $\alpha_i^p-\alpha_j$, divisible by a large power of $p$, while it is small when $M(P)$ is close to $1$.
\end{remark}

\section{The Jacobian conjecture}

A polynomial map $F=(F_1,\dots,F_n):k^n\to k^n$ has a \emph{Jacobian matrix} $J_F=\bigl(\frac{\partial F_i}{\partial x_j}\bigr)$, with polynomial
entries (formal derivatives). If $F$ has a polynomial inverse $G$, the chain rule applied to $G\circ F=\mathrm{id}$ gives $J_G(F(x))J_F(x)=I$, so
$\det J_F$ is a polynomial with a polynomial inverse: a non-zero constant (Exercise~\ref{conjectures:ex:jacdim1}). In 1939 Ott-Heinrich Keller asked for the converse.

\begin{center}
\emph{The Jacobian conjecture.} Let $F:\C^n\to\C^n$ be a polynomial map with $\det J_F$ a non-zero constant. Then $F$ is bijective and $F^{-1}$ is a
polynomial map.
\end{center}

It is true for $n=1$ (Exercise~\ref{conjectures:ex:jacdim1}) and open for every $n\ge2$. Over a field of characteristic $p$ it is false: $x\mapsto x-x^p$ has derivative
$1$ but sends $0$ and $1$ to $0$ (Exercise~\ref{conjectures:ex:charpjac}). The conjecture has collected an unusual number of incorrect proofs, some of them published;
the survey of Bass, Connell and Wright (1982) tells the early history.

\begin{think}
For $F(x,y)=(x+y^2,\,y)$, compute $\det J_F$ and the inverse. Then do the same for $F(x,y)=(x+y^2,\,y+(x+y^2)^2)$. What is the degree of the inverse?
\end{think}

For maps of degree $2$ there is a two-line identity.

\begin{theorem}[Wang, 1980]\label{conjectures:thm:wang}
Let $k$ be a field of characteristic $\ne2$ and $F:k^n\to k^n$ a polynomial map whose components have degree at most $2$. If $\det J_F(x)\ne0$ for every
$x\in k^n$, then $F$ is injective on $k^n$.
\end{theorem}
\begin{proof}
We show that for all $a,b\in k^n$,
\[
F(b)-F(a)=J_F\Bigl(\frac{a+b}2\Bigr)(b-a).
\]
It is enough to check it for one component $F_i=c+\ell+Q$, with $c$ a constant, $\ell$ linear and $Q$ a quadratic form. Since $2$ is invertible in $k$, we can
write $Q(x)=x^{\mathsf T}Ax$ with $A$ symmetric (put half of the coefficient of $x_jx_l$ in each of the entries $(j,l)$ and $(l,j)$). The gradient of $F_i$
at $m$ is $\ell+2Am$ (as a row vector, $\ell^{\mathsf T}+2m^{\mathsf T}A$), so at $m=\frac{a+b}2$ the right-hand side is $\ell(b-a)+(a+b)^{\mathsf T}A(b-a)$.
On the left, $F_i(b)-F_i(a)=\ell(b-a)+b^{\mathsf T}Ab-a^{\mathsf T}Aa$, and expanding with $a^{\mathsf T}Ab=b^{\mathsf T}Aa$,
\[
(a+b)^{\mathsf T}A(b-a)=a^{\mathsf T}Ab-a^{\mathsf T}Aa+b^{\mathsf T}Ab-b^{\mathsf T}Aa=b^{\mathsf T}Ab-a^{\mathsf T}Aa.
\]
Now if $F(a)=F(b)$, the identity gives $J_F(\frac{a+b}2)(b-a)=0$, and since $J_F(\frac{a+b}2)$ is invertible, $b=a$.
\end{proof}

Over $\R$ and $\C$ the identity is the mean value theorem with the exact midpoint: for a quadratic map the derivative is affine, and the average of an affine
function over a segment is its value at the midpoint. Injectivity is half of the conjecture. For the other half we turn to Ax and Grothendieck.

\begin{remark}[what is known]\label{conjectures:rem:jacknown}
Wang proved more: over $\C$, a quadratic map with constant non-zero Jacobian has a polynomial inverse. Bass, Connell and Wright (1982), and independently Yagzhev (1980), showed
that it suffices to prove the conjecture, in all dimensions, for the maps $x+H(x)$ with $H$ homogeneous of degree $3$; Dru\.zkowski (1983) reduced further to
$H(x)=\bigl((\lambda_1\cdot x)^3,\dots,(\lambda_n\cdot x)^3\bigr)$. So degree $2$ is proved and degree $3$ is the whole problem. In dimension $2$, Moh (1983)
proved the conjecture for maps of degree at most $100$. When $F^{-1}$ exists it has degree at most $(\deg F)^{n-1}$ (Bass--Connell--Wright, after Gabber).
Over $\R$, the \emph{strong real Jacobian conjecture} (if $\det J_F$ never vanishes on $\R^n$, then $F$ is injective) was shown false by Pinchuk
(1994), with a map of the plane of high degree; Theorem~\ref{conjectures:thm:wang} shows that degree $2$ is too small for such a counterexample.
\end{remark}

\section{Injective implies surjective}

A map from a finite set to itself is injective if and only if it is surjective. Ax and Grothendieck turned this triviality into a theorem about $\C$.

\begin{theorem}\label{conjectures:thm:axgr}
Let $K$ be an algebraic extension of $\mathbb F_p$, for example an algebraic closure, and $F:K^n\to K^n$ a polynomial map with coefficients in $K$. If $F$ is
injective, then $F$ is surjective.
\end{theorem}
\begin{proof}
Let $b\in K^n$. Let $k_0\subseteq K$ be the subfield generated over $\mathbb F_p$ by the finitely many coefficients of $F$ and the coordinates $b_1,\dots,b_n$.
Each generator is algebraic over $\mathbb F_p$, so $k_0$ has finite degree over $\mathbb F_p$ by the tower law (Tome~2, Chapter~8): it is a finite field.
Since the coefficients of $F$ lie in $k_0$, $F$ maps the finite set $k_0^n$ to itself, injectively. So it is a bijection of $k_0^n$, and $b\in k_0^n$ has a
preimage.
\end{proof}

\begin{corollary}\label{conjectures:cor:wangfin}
Let $K$ be an algebraic extension of $\mathbb F_p$, $p$ odd, and $F:K^n\to K^n$ a polynomial map of degree at most $2$ with $\det J_F(x)\ne0$ for all $x\in K^n$.
Then $F$ is a bijection of $K^n$.
\end{corollary}
\begin{proof}
Theorem~\ref{conjectures:thm:wang} gives injectivity and Theorem~\ref{conjectures:thm:axgr} surjectivity.
\end{proof}

\begin{remark}[the transfer to $\C$]\label{conjectures:rem:axC}
The theorem holds over $\C$: an injective polynomial map $\C^n\to\C^n$ is surjective (Ax, 1968; Grothendieck, 1966, by a different route). Ax's proof
transfers Theorem~\ref{conjectures:thm:axgr} to every algebraically closed field of characteristic $0$. The statement ``every injective polynomial map of degree at most
$d$ in $n$ variables is surjective'' is a single first-order sentence in the language of fields, true in the algebraic closure of every $\mathbb F_p$; a theorem
of model theory (the compactness theorem, with the completeness of the theory of algebraically closed fields of a given characteristic) says that such a
sentence, true in characteristic $p$ for all large $p$, is true in $\C$. A purely algebraic version uses the Nullstellensatz to reduce a counterexample
over $\C$ to one over a finite field. Combined with Theorem~\ref{conjectures:thm:wang}, it shows that a quadratic map $\C^n\to\C^n$ with non-zero constant
Jacobian is bijective; by a theorem of Bia\l ynicki-Birula and Rosenlicht (1962) the inverse of a polynomial bijection of $\C^n$ is polynomial.
\end{remark}

Exercise~\ref{conjectures:ex:axfails} shows what goes wrong over $\Q$ and over $\mathbb F_p(t)$, where the fields are not unions of finite (or algebraically closed)
pieces.

\section{Why it is hard}

The three open problems resist for three different reasons. For the inverse Galois problem, a group must be \emph{built}: every known method produces an
extension with a prescribed group from geometry (covers of the line, rigid triples of conjugacy classes), and each family of groups needs its own geometry;
there is no general construction, and no known obstruction either. For Lehmer's problem, the proof of Theorem~\ref{conjectures:thm:smyth} uses that a non-reciprocal
$P$ differs from $P^*$ in some coefficient, an integer gap of size at least $1$; for reciprocal polynomials $P/P^*=\pm1$ and the argument has nothing to hold
on to. Dobrowolski's method, which works for all polynomials, loses a power of $\log d$ that nobody knows how to recover. For the Jacobian conjecture, the
hypothesis is local (the derivative is invertible at every point) and the conclusion global; Theorem~\ref{conjectures:thm:wang} bridges the gap only because a quadratic
map is determined by its derivative at a single midpoint, and the reduction to cubic maps shows that the first non-trivial case already contains everything.

\section{Reading a paper: Smyth on non-reciprocal algebraic integers}

C.~J.~Smyth's \emph{On the product of the conjugates outside the unit circle of an algebraic integer} (Bulletin of the London Mathematical Society, 1971) is
seven pages long. Read the introduction for the statement and the equality case. The proof begins exactly as Theorem~\ref{conjectures:thm:smyth}: the rational function
$P/P^*$ has integer coefficients and splits into two Blaschke-type products bounded by $1$ in the disc. Where we used only $|\gamma_k|\le1-|\gamma_0|^2$,
Smyth extracts more from Schur's description of the coefficients of bounded functions; compare his inequalities with
Lemma~\ref{conjectures:lem:coef} and see where the gap between $1.2808$ and $1.3247$ closes. The idea worth stealing: \emph{an integer that is not zero has absolute
value at least $1$}, applied to a power series rather than to a number. The same idea, with a congruence modulo $m$ instead, gives Exercise~\ref{conjectures:ex:oddcoef}.

For the Jacobian conjecture, read Sections~I and~II of Bass, Connell and Wright, \emph{The Jacobian conjecture: reduction of degree and formal expansion
of the inverse} (Bulletin of the AMS, 1982): the history, the elementary facts of Exercise~\ref{conjectures:ex:jacdim1}, and the reduction to cubic maps, whose trick is
to add variables to lower the degree.

\section{Open doors}

\begin{itemize}
\item Extend Figure~\ref{conjectures:fig:mahler} to degree $20$ with a smarter search (only reciprocal polynomials, roots counted with Sturm sequences). At which degrees
do new Salem numbers below $1.3$ appear?
\item For each transitive subgroup $G$ of $S_5$, find a polynomial of small height with group $G$, and estimate from the Frobenius cycle types (Tome~12,
Chapter~8) how often each group occurs among random quintics of height $\le H$.
\item Test the Jacobian conjecture on random cubic maps $x+H(x)$ over $\mathbb F_p$ with $\det J=1$: are they always injective on $\mathbb F_p^n$? Find the
smallest $p$ and $n$ for which you can search all such maps.
\item Lehmer's problem for polynomials with coefficients in $\{-1,1\}$ (the Littlewood polynomials): what is the smallest measure $>1$ you can find in
degree up to $30$, and is it attained by a reciprocal polynomial?
\end{itemize}

\section{Exercises}

\begin{exercise}[$\star\star$]\label{conjectures:ex:period}
Let $\zeta=e^{2i\pi/11}$ and $\eta=\zeta+\zeta^{-1}=2\cos\frac{2\pi}{11}$. Show that $\Q(\eta)$ is a Galois extension of $\Q$ with group $\Z/5\Z$, and that
the minimal polynomial of $\eta$ is $X^5+X^4-4X^3-3X^2+3X+1$.
\end{exercise}

\begin{exercise}[$\star\star$]\label{conjectures:ex:x4}
Show that the Galois group of $X^4-X-1$ over $\Q$ is $S_4$, using its factorisations modulo $2$, $7$ and $17$:
$X^4+X+1$ (irreducible over $\mathbb F_2$), $(X-3)(X^3+3X^2+2X-2)$ over $\mathbb F_7$, and $(X+2)(X+5)(X^2-7X+5)$ over $\mathbb F_{17}$. Its discriminant
is $-283$.
\end{exercise}

\begin{exercise}[$\star\star$]\label{conjectures:ex:dihedral}
Show that the symmetries of a square, acting on its four vertices, form a transitive subgroup of $S_4$ of order $8$ that contains a $4$-cycle and a
transposition. Why does this not contradict Lemma~\ref{conjectures:lem:sngen}, nor the theorem that a transitive subgroup of $S_p$ containing a transposition is $S_p$
for $p$ prime?
\end{exercise}

\begin{exercise}[$\star$]\label{conjectures:ex:mahlerbasic}
(a) Show that $M(P^*)=M(P)$ when $P(0)\ne0$, and that $M(P(X^k))=M(P)$ for $k\ge1$. (b) Show that $X^3-X-1$ is irreducible and not reciprocal, and that
$M(X^3-X-1)$ is its real root $\theta_0$. Deduce that the best constant $c$ in Corollary~\ref{conjectures:cor:nonrec} satisfies $1.2807\le c\le1.3248$.
\end{exercise}

\begin{exercise}[$\star\star\star$]\label{conjectures:ex:oddcoef}
Let $P\in\Z[X]$ be monic, irreducible and not reciprocal, with all coefficients odd. Show that $M(P)\ge\frac{1+\sqrt5}2$. (This is a theorem of Borwein,
Hare and Mossinghoff, 2004; $X^2-X-1$ shows that it is sharp.)
\end{exercise}

\begin{exercise}[$\star$]\label{conjectures:ex:jacdim1}
(a) Show that if a polynomial map $F:\C^n\to\C^n$ has a polynomial inverse, then $\det J_F$ is a non-zero constant. (b) Prove the Jacobian conjecture
for $n=1$.
\end{exercise}

\begin{exercise}[$\star\star$]\label{conjectures:ex:charpjac}
(a) Show that $x\mapsto x-x^p$ has derivative $1$ on $\mathbb F_p$ and on $\overline{\mathbb F}_p$, and is neither injective nor, on $\mathbb F_p$, surjective.
Is it surjective on an algebraic closure? (b) Show that Theorem~\ref{conjectures:thm:wang} fails in characteristic $2$.
\end{exercise}

\begin{exercise}[$\star$]\label{conjectures:ex:axfails}
Show that $x\mapsto x^3$ is an injective polynomial map $\Q\to\Q$ which is not surjective, and that $x\mapsto x^p$ is an injective polynomial map of
$\mathbb F_p(t)$ to itself which is not surjective. Where does the proof of Theorem~\ref{conjectures:thm:axgr} break down?
\end{exercise}

\begin{exercise}[$\star\star$]\label{conjectures:ex:triangular}
Let $F(x,y)=\bigl(x+y^2,\ y+(x+y^2)^2\bigr)$. Compute $\det J_F$ directly, then find $F^{-1}$ by writing $F$ as a composition of two maps of the form
$(x,y)\mapsto(x+g(y),y)$ and $(x,y)\mapsto(x,y+g(x))$. What are the degrees of $F$ and $F^{-1}$?
\end{exercise}

\section{Hints}

\begin{hint}{conjectures:ex:period}
$\eta$ is fixed exactly by the subgroup $\{\pm1\}$ of $(\Z/11\Z)^*$, which is cyclic of order $10$. The conjugates of $\eta$ are $2\cos\frac{2\pi k}{11}$,
$k=1,\dots,5$; use $\zeta^{11}=1$ and $1+\zeta+\dots+\zeta^{10}=0$ to compute their elementary symmetric functions.
\end{hint}
\begin{hint}{conjectures:ex:x4}
Check that $283$ is prime, then read off a $4$-cycle, a $3$-cycle and a transposition.
\end{hint}
\begin{hint}{conjectures:ex:dihedral}
Number the vertices $1,2,3,4$ around the square. The reflection in a diagonal is a transposition.
\end{hint}
\begin{hint}{conjectures:ex:mahlerbasic}
(a) The roots of $P^*$ are the $1/\alpha_i$, and its leading coefficient is $P(0)$. (b) A reducible cubic has a rational root, which must be $\pm1$.
\end{hint}
\begin{hint}{conjectures:ex:oddcoef}
Follow the proof of Theorem~\ref{conjectures:thm:smyth}, and show that $q_k$ is even: modulo $2$, $P\equiv P^*$.
\end{hint}
\begin{hint}{conjectures:ex:jacdim1}
(a) Chain rule; a polynomial with a polynomial inverse has degree $0$. (b) $\deg F'=\deg F-1$.
\end{hint}
\begin{hint}{conjectures:ex:charpjac}
(a) The roots of $x^p-x-c$ in an algebraic closure. (b) $x\mapsto x+x^2$.
\end{hint}
\begin{hint}{conjectures:ex:axfails}
Is $\Q$ a union of finite subfields? Is $t$ a $p$-th power in $\mathbb F_p(t)$?
\end{hint}
\begin{hint}{conjectures:ex:triangular}
$F=H\circ G$ with $G(x,y)=(x+y^2,y)$ and $H(u,v)=(u,v+u^2)$.
\end{hint}

\section{Solutions}

\begin{solution}{conjectures:ex:period}
\idea Galois theory turns ``which field'' into ``which subgroup'', and the relation $1+\zeta+\dots+\zeta^{10}=0$, rewritten in $\eta$, gives the polynomial.
\writeup By Tome~2, Exercise~10.8,
\[
\operatorname{Gal}(\Q(\zeta)/\Q)\cong(\Z/11\Z)^*,
\]
with $a$ acting by $\sigma_a(\zeta)=\zeta^a$; this group is cyclic of order $10$ (Tome~1, Theorem~5.3). Since $\sigma_a(\eta)=\zeta^a+\zeta^{-a}=2\cos\frac{2\pi a}{11}$, we have $\sigma_a(\eta)=\eta$ if and only if
$a\equiv\pm1\pmod{11}$. So $\Q(\eta)$ is the fixed field of $\{\pm1\}$, a normal subgroup since the group is abelian, and by Tome~2, Theorem~10.6, it is
Galois over $\Q$ with group $(\Z/11\Z)^*/\{\pm1\}$, cyclic of order $5$; in particular $[\Q(\eta):\Q]=5$. Now divide $1+\zeta+\dots+\zeta^{10}=0$ by
$\zeta^5$: $1+\sum_{k=1}^5(\zeta^k+\zeta^{-k})=0$. With $\zeta^k+\zeta^{-k}=\eta(\zeta^{k-1}+\zeta^{1-k})-(\zeta^{k-2}+\zeta^{2-k})$,
\[
\begin{aligned}
\zeta^2+\zeta^{-2}&=\eta^2-2,\\
\zeta^3+\zeta^{-3}&=\eta^3-3\eta,\\
\zeta^4+\zeta^{-4}&=\eta^4-4\eta^2+2,\\
\zeta^5+\zeta^{-5}&=\eta^5-5\eta^3+5\eta,
\end{aligned}
\]
and the sum is $\eta^5+\eta^4-4\eta^3-3\eta^2+3\eta+1=0$. A monic polynomial of degree $5=[\Q(\eta):\Q]$ vanishing at $\eta$ is its minimal polynomial.
\end{solution}

\begin{solution}{conjectures:ex:x4}
\idea Dedekind's theorem turns each factorisation into a cycle type; we need three of them.
\writeup $283$ is prime (not divisible by $2,3,5,7,11,13$ and $17^2>283$), so $2$, $7$, $17$ do not divide the discriminant. Modulo $2$ the polynomial is
irreducible, so $f$ is irreducible over $\Q$ (as in the proof of Theorem~\ref{conjectures:thm:sn}) and $G$ is transitive. By Tome~12, Theorem~2.20, $G$ contains a
permutation of type $(1,3)$ (modulo $7$), a $3$-cycle, and one of type $(1,1,2)$ (modulo $17$), a transposition. Lemma~\ref{conjectures:lem:sngen} with $n=4$ gives
$G=S_4$. The factorisations can be checked by expanding: for instance $3^4-3-1=77=7\cdot11$.
\end{solution}

\begin{solution}{conjectures:ex:dihedral}
\idea The lemma needs an $(n-1)$-cycle, and $n=4$ is not prime.
\writeup With the vertices $1,2,3,4$ in order, the rotations are $\mathrm{id}$, $(1\,2\,3\,4)$, $(1\,3)(2\,4)$, $(1\,4\,3\,2)$, and the reflections are
$(1\,3)$, $(2\,4)$ (diagonals), $(1\,2)(3\,4)$, $(1\,4)(2\,3)$. The group is transitive (the rotations already are), has order $8\ne24$, contains the
$4$-cycle $(1\,2\,3\,4)$ and the transposition $(1\,3)$. It contains no $3$-cycle, so Lemma~\ref{conjectures:lem:sngen} does not apply. The theorem for $S_p$ uses that
a $p$-cycle generates a group of prime order, all of whose non-trivial elements are $p$-cycles; for $n=4$, the square $(1\,3)(2\,4)$ of the $4$-cycle is
not a $4$-cycle, and that is where the argument breaks.
\end{solution}

\begin{solution}{conjectures:ex:mahlerbasic}
\idea Track the roots.
\writeup (a) If $P=a\prod(X-\alpha_i)$ with all $\alpha_i\ne0$, then
\[
P^*=a\prod(1-\alpha_iX)=a\prod(-\alpha_i)\prod(X-\alpha_i^{-1}),
\]
so
\[
M(P^*)=|a|\prod|\alpha_i|\prod\max(1,|\alpha_i|^{-1})=|a|\prod\max(|\alpha_i|,1)=M(P).
\]
The roots of $P(X^k)$ are the $k$-th roots $\beta$ of the $\alpha_i$, $k$ for each, with $|\beta|^k=|\alpha_i|$; so
\[
\prod_\beta\max(1,|\beta|)=\prod_i\max\bigl(1,|\alpha_i|^{1/k}\bigr)^k=\prod_i\max(1,|\alpha_i|).
\]
(b) $\pm1$ are not roots ($1-1-1=-1$, $-1+1-1=-1$), so the cubic has no rational root and is irreducible. Its reverse is $-X^3-X^2+1\ne\pm(X^3-X-1)$. It has
one real root $\theta_0\approx1.3247$ ($f(1)<0<f(2)$, and $f'=3X^2-1$ shows the local maximum $f(-1/\sqrt3)<0$), and the complex roots have product
$1/\theta_0$ in modulus, each of modulus $\theta_0^{-1/2}<1$. So $M=\theta_0$. The best constant $c$ is at least the bound $\frac{1+\sqrt{17}}4>1.2807$ of
Corollary~\ref{conjectures:cor:nonrec}, and at most $\theta_0<1.3248$ because of this example (Smyth's theorem, Remark~\ref{conjectures:rem:smythsharp}, says $c=\theta_0$).
\end{solution}

\begin{solution}{conjectures:ex:oddcoef}
\idea The same proof, with an integer gap of $2$ instead of $1$.
\writeup $|P(0)|$ is odd; if it is at least $3$, $M(P)\ge3$. Otherwise $P(0)=\pm1$, and the proof of Theorem~\ref{conjectures:thm:smyth} applies word for word up to
the last line. Modulo $2$, every coefficient of $P$ and of $P^*$ is $1$, so $P\equiv P^*$ in $\mathbb F_2[[z]]$. Since $P=GP^*$ and $P^*$ is invertible in
$\mathbb F_2[[z]]$ (constant term $1$), $G\equiv1\pmod2$; so $q_k$, the first non-zero coefficient after $q_0$, is even, and $|q_k|\ge2$. The final inequality becomes
$\frac2M\le2\bigl(1-\frac1{M^2}\bigr)$, that is $M^2-M-1\ge0$, $M\ge\frac{1+\sqrt5}2$. For $X^2-X-1$, the roots are $\frac{1\pm\sqrt5}2$ and
$M=\frac{1+\sqrt5}2$; it is not reciprocal since its reverse is $-X^2-X+1$.
\end{solution}

\begin{solution}{conjectures:ex:jacdim1}
\idea Derivatives of compositions multiply.
\writeup (a) If $G\circ F=\mathrm{id}$ with $G$ polynomial, the chain rule (valid for formal derivatives) gives $J_G(F(x))J_F(x)=I$, so
$\det J_G(F(x))\cdot\det J_F(x)=1$ as polynomials. A product of two polynomials equal to $1$ has degree $0$ in both, so $\det J_F$ is a non-zero constant.
(b) If $F'=c\ne0$ is constant, then $F=cx+d$, which has the inverse $y\mapsto(y-d)/c$.
\end{solution}

\begin{solution}{conjectures:ex:charpjac}
\idea In characteristic $p$, the derivative of $x^p$ is $px^{p-1}=0$.
\writeup (a) $\frac{\dd}{\dd x}(x-x^p)=1-px^{p-1}=1$. By Fermat, $x-x^p=0$ for every $x\in\mathbb F_p$: the map is constant on $\mathbb F_p$, so neither injective
nor surjective ($p\ge2$). On an algebraic closure it is surjective, since $x^p-x+c$ has a root for every $c$, but not injective ($0,1\mapsto0$); this does
not contradict Theorem~\ref{conjectures:thm:axgr}, which goes from injective to surjective, not back. (b) Over $\mathbb F_2$, $F(x)=x+x^2$ has $F'=1+2x=1$, never zero,
and $F(0)=F(1)=0$. The proof fails because $Q(x)=x^2$ cannot be written $x^{\mathsf T}Ax$ with the factor $2$ in the gradient.
\end{solution}

\begin{solution}{conjectures:ex:axfails}
\idea The proof needs every finite piece of the problem to live in a finite field.
\writeup $x^3=y^3$ in $\Q$ implies $x=y$ (cube is strictly increasing on $\R$), and $2$ is not the cube of a rational ($\sqrt[3]2$ is irrational). In
$\mathbb F_p(t)$, $x^p=y^p$ gives $(x-y)^p=0$, so $x=y$; and $t$ is not a $p$-th power: if $t=(a/b)^p$ with polynomials $a,b$, then $tb^p=a^p$, and the
degree of the left side is $\equiv1$, of the right side $\equiv0$ modulo $p$. The proof of Theorem~\ref{conjectures:thm:axgr} takes the subfield generated by the data:
in $\Q$ it is $\Q$ itself, and in $\mathbb F_p(t)$ it contains $t$; neither is finite.
\end{solution}

\begin{solution}{conjectures:ex:triangular}
\idea Triangular maps are inverted by subtraction, and compositions are inverted in the reverse order.
\writeup With $u=x+y^2$, $F=(u,y+u^2)$, and
\[
J_F=\begin{pmatrix}1&2y\\2u&1+4uy\end{pmatrix},\qquad\det J_F=1+4uy-4uy=1.
\]
$G(x,y)=(x+y^2,y)$ has inverse $G^{-1}(u,v)=(u-v^2,v)$, and $H(u,v)=(u,v+u^2)$ has inverse $H^{-1}(u,v)=(u,v-u^2)$. Since $F=H\circ G$,
\[
F^{-1}(u,v)=G^{-1}\bigl(H^{-1}(u,v)\bigr)=\bigl(u-(v-u^2)^2,\ v-u^2\bigr).
\]
$F$ has degree $4$ and so does $F^{-1}$, within the bound $(\deg F)^{n-1}=4$ of Remark~\ref{conjectures:rem:jacknown}.
\end{solution}

\section*{Sources and further reading}

\begin{itemize}
\item D.~Hilbert, \emph{\"Uber die Irreducibilit\"at ganzer rationaler Functionen mit ganzzahligen Coefficienten}, Journal f\"ur die reine und angewandte
Mathematik 110 (1892), 104--129. E.~Noether, \emph{Gleichungen mit vorgeschriebener Gruppe}, Mathematische Annalen 78 (1918), 221--229. R.~G.~Swan,
\emph{Invariant rational functions and a problem of Steenrod}, Inventiones Mathematicae 7 (1969), 148--158.
\item G.~Malle and B.~H.~Matzat, \emph{Inverse Galois Theory}, second edition, Springer, 2018, for the rigidity method and the state of the art; and
J.-P.~Serre, \emph{Topics in Galois Theory}, second edition, A~K~Peters, 2008. J.~G.~Thompson, \emph{Some finite groups which appear as $\operatorname{Gal}L/K$,
where $K\subseteq\Q(\mu_n)$}, Journal of Algebra 89 (1984), 437--499. On Shafarevich's theorem and its repair: J.~Neukirch, A.~Schmidt and K.~Wingberg,
\emph{Cohomology of Number Fields}, second edition, Springer, 2008, Chapter~IX.
\item X.~Huang, B.~Jackson, K.-H.~Lee, B.~Poonen, R.~Pries and S.~Zhang, \emph{The Mathieu group $M_{23}$ is a Galois group over $\Q$}, preprint,
arXiv:2608.08538 (2026).
\item B.~L.~van der Waerden, \emph{Die Seltenheit der reduziblen Gleichungen und der Gleichungen mit Affekt}, Monatshefte f\"ur Mathematik und Physik 43
(1936). M.~Bhargava, \emph{Galois groups of random integer polynomials and van der Waerden's conjecture}, Annals of Mathematics 201 (2025).
H.~Osada, \emph{The Galois groups of the polynomials $X^n+aX^l+b$}, Journal of Number Theory 25 (1987), 230--238.
\item D.~H.~Lehmer, \emph{Factorization of certain cyclotomic functions}, Annals of Mathematics 34 (1933), 461--479. C.~J.~Smyth, \emph{On the product of
the conjugates outside the unit circle of an algebraic integer}, Bulletin of the London Mathematical Society 3 (1971), 169--175. C.~Smyth,
\emph{The Mahler measure of algebraic numbers: a survey}, in \emph{Number Theory and Polynomials}, London Mathematical Society Lecture Note Series 352,
Cambridge, 2008, 322--349, covers the whole subject, including Breusch's 1951 bound and Smyth's theorem with its equality case. E.~Dobrowolski, \emph{On a
question of Lehmer and the number of irreducible factors of a polynomial}, Acta Arithmetica 34 (1979), 391--401.
\item M.~J.~Mossinghoff, G.~Rhin and Q.~Wu, \emph{Minimal Mahler measures}, Experimental Mathematics 17 (2008), 451--458. P.~Borwein, K.~G.~Hare and
M.~J.~Mossinghoff, \emph{The Mahler measure of polynomials with odd coefficients}, Bulletin of the London Mathematical Society 36 (2004), 332--338.
\item O.-H.~Keller, \emph{Ganze Cremona-Trans\-for\-ma\-tio\-nen}, Monatshefte f\"ur Mathematik und Physik 47 (1939), 299--306. H.~Bass, E.~H.~Connell and
D.~Wright, \emph{The Jacobian conjecture: reduction of degree and formal expansion of the inverse}, Bulletin of the AMS 7 (1982), 287--330. S.~S.-S.~Wang,
\emph{A Jacobian criterion for separability}, Journal of Algebra 65 (1980), 453--494. A.~van den Essen, \emph{Polynomial Automorphisms and the Jacobian
Conjecture}, Birkh\"auser, 2000, for Yagzhev's and Dru\.zkowski's reductions, Moh's theorem and Pinchuk's counterexample (S.~Pinchuk,
\emph{A counterexample to the strong real Jacobian conjecture}, Mathematische Zeitschrift 217 (1994), 1--4).
\item J.~Ax, \emph{The elementary theory of finite fields}, Annals of Mathematics 88 (1968), 239--271. A.~Grothendieck, \emph{\'El\'ements de g\'eom\'etrie
alg\'ebrique IV}, troisi\`eme partie, Publications Math\'ematiques de l'IH\'ES 28 (1966), Proposition~10.4.11. A.~Bia\l ynicki-Birula and M.~Rosenlicht,
\emph{Injective morphisms of real algebraic varieties}, Proceedings of the AMS 13 (1962), 200--203. T.~Tao's blog post \emph{Infinite fields, finite fields, and the
Ax--Grothendieck theorem} (2009) gives the transfer to $\C$ in elementary language.
\item The polynomials of Theorem~\ref{conjectures:thm:sn} and the data of Figures~\ref{conjectures:fig:lehmer} and~\ref{conjectures:fig:mahler} were computed for this book.
\end{itemize}
